\documentclass[11pt]{article}
\usepackage[margin=1in]{geometry}
\usepackage{amsmath,amssymb}
\usepackage[numbers]{natbib}
\usepackage{times}
\usepackage[hidelinks]{hyperref}
\usepackage{enumitem}
\usepackage{titlesec}
\titleformat{\section}{\large\bfseries}{\thesection}{1em}{}
\titleformat{\subsection}{\normalsize\bfseries}{\thesubsection}{1em}{}

\title{Why No One Talks Backwards: \\
Developmental Exploration, Reachability, and the Absence of \\
Temporal Reversal as a Grammatical Operation}
\author{Flyxion \\ \small Independent Researcher}
\date{}

\begin{document}
\maketitle

\begin{abstract}
\noindent Natural languages recruit an enormous variety of formal operations for grammatical purposes: reduplication, ablaut, consonant mutation, tone, stress shift, infixation, metathesis, subtraction, and suppletion, among others. I am unaware of any attested natural language in which acoustic time-reversal of a form productively realizes a grammatical category such as tense, aspect, number, negation, or evidentiality. This absence is not explained by the physical impossibility of backward speech, since humans can and do learn to produce an approximation of it under unusual conditions. Nor is it explained by a general claim that human phonology is limited to sounds infants happen to produce during babbling, since mature speakers routinely master articulations well outside their own infantile repertoire. The explanation offered here instead turns on a distinction between three properties that are easily conflated: physical possibility, learnability once demonstrated, and developmental or cultural \emph{discoverability}, meaning the likelihood that a transformation will ever be generated as a candidate by the ordinary processes through which languages are produced, transmitted, and reconstructed across generations. A further distinction, between acoustic reversal of a waveform and segmental reversal of an already-discretized phonological representation, separates the transformation actually at issue from the reordering operations, such as metathesis, that languages plainly do exhibit. Drawing on developmental psychology, the cross-linguistic babbling literature, the empirical literature on fluent backward talkers, iterated-learning experiments in cultural evolution, and a general theory of variation-before-selection, the essay argues that grammatical operations are restricted to transformations that fall within the reachable neighborhood of ordinary articulatory and perceptual variation, and that acoustic reversal of a signal does not correspond, in general, to reversal of the articulatory gestures that produced it, so that it never appears as a naturally occurring perturbation for cultural selection to act upon. A concluding section offers, as an explicitly separate and more conjectural claim, a generalization of the same reachability principle: a 2026 multi-agent AI coordination incident, in which isolated agents turned an ordinary shared cache into an unintended communication channel, illustrates the mirror-image case, in which a transformation nobody designed nonetheless proves locally reachable and is retained.
\end{abstract}

\section{The Question, Narrowed}

The question in the title admits several readings, and most of them are the wrong ones. It is not a claim that human beings cannot produce reversed speech: expert backward talkers exist, and their skill has been studied directly \citep{cowan1982,cowan1985}. It is not a claim that natural phonological inventories are restricted to sounds that infants happen to stumble upon while babbling, since adult speakers routinely acquire, through instruction and practice, articulations that do not appear spontaneously in infant vocalization. Languages themselves already contain reversal-like operations at the segmental level, most notably metathesis, and language games in many cultures perform quite elaborate reorderings of a word's material. And it is not a claim about writing systems or reading direction, since the direction in which a script is transcribed onto a page has no necessary relationship to the temporal structure of the speech it represents.

The narrower and more defensible question requires separating three operations that ordinary usage of the word ``reversal'' blurs together. Given an acoustic signal $x(t)$ for $0 \le t \le T$, define acoustic time-reversal as
\[
R_a[x](t) \;=\; x(T-t),
\]
the literal reversal of the waveform itself. Given a phonological representation as a sequence of discrete segments $(\phi_1, \phi_2, \ldots, \phi_n)$, define segmental reversal as
\[
R_p(\phi_1, \phi_2, \ldots, \phi_n) \;=\; (\phi_n, \ldots, \phi_2, \phi_1),
\]
the reordering of an already-segmented symbolic representation. And define metathesis as the transposition of a small number of adjacent or near-adjacent segments within an otherwise unchanged sequence, a local special case of $R_p$ rather than a global one. These three operations are not interchangeable, and, as later sections argue, treating them as equivalent is precisely what makes the question look either trivially false (languages obviously have metathesis) or paradoxically strong (no language reverses its acoustic signal).

With that distinction in place, the claim this essay defends is narrow and, I hope, falsifiable: I am unaware of any attested natural language in which $R_a$, acoustic time-reversal of a form, productively realizes a grammatical category such as tense, aspect, number, negation, or evidentiality. Segmental reversal $R_p$ and local metathesis are attested, in various languages, as historical sound changes and, occasionally, as marginal morphological processes; global acoustic reversal, so far as I am aware, is not attested anywhere as an ordinary grammatical device.

This is a different question from why people do not casually speak in reverse for amusement, since some people do exactly that, given sufficient practice and, increasingly, given access to recording technology that lets them hear a target they could not otherwise construct. Separating these two questions from one another is the analytic move on which the rest of the essay depends.

\section{Explore and Exploit in Development}

A useful starting point is an asymmetry that developmental psychologists have long noted between the behavior of young children and that of adults. An adult who wants coffee does not experimentally combine substances in the kitchen to discover which ones produce an acceptable beverage; adults act on an accumulated causal model of their environment and select actions expected to produce known outcomes. Alison Gopnik and colleagues have argued that childhood occupies a structurally different position in what is, in reinforcement-learning terms, an explore-exploit tradeoff: children tolerate, and in fact generate, extraordinary quantities of behavior with no evident payoff, because they are acquiring a causal model of a world they do not yet understand \citep{gopnik1999,gopnik2004,gopnik2020}. Where an exploiting agent asks which action has the highest expected value under its current model, an exploring agent sometimes selects an action precisely because its outcome is uncertain: the informational value of finding out is itself a payoff.

This reframes a piece of everyday interaction that is usually treated as unremarkable. An adult asking a child \emph{why did you do that} is presupposing a particular causal grammar for behavior, roughly \emph{goal $\rightarrow$ deliberation $\rightarrow$ action}, in which a legitimate action must be traceable to an antecedent instrumental objective. A great deal of childhood behavior, however, is better modeled as \emph{uncertainty $\rightarrow$ intervention $\rightarrow$ observation $\rightarrow$ updated model}. Under that model, \emph{to see what would happen} is not an evasive non-answer; it is a reasonably precise description of the causal structure of the action, and it is a recognizable precursor to the controlled experimentation that the scientific method later formalizes, though it should not be identified with science outright: indiscriminate intervention lacks the systematic control, explicit hypothesis, measurement, and replication that distinguish an experiment from an exploratory perturbation. The apparent absurdity of the child's answer says more about the adult's implicit theory of action than about any deficiency in the child's reasoning.

Vocal and gestural babbling are the clearest illustrations of this exploratory regime, and they are also the direct antecedents of the phenomenon this essay is about.

\section{Babble as Exploration of a Motor Possibility Space}

An infant's vocal tract affords an enormous space of possible configurations, and the infant does not begin with anything resembling a phonological inventory and proceed to execute its members. It moves through the space. Canonical babbling research documents a broad early period of vocal production, only later shaped by the ambient language into an increasingly narrow, increasingly reproducible set of categories \citep{oller2000}. Patricia Kuhl's work on early phonetic perception documents a parallel narrowing on the perceptual side, as infants who initially discriminate the full range of human phonetic contrasts become progressively tuned to the contrasts of their linguistic environment \citep{kuhl2004}. De Boysson-Bardies and colleagues' cross-linguistic work gives a more qualified version of this concurrency than is sometimes claimed. Their 1984 study found that adult listeners could identify a French, Arabic, or Cantonese infant's ambient language from recordings of ten-month-old canonical babbling above chance, using primarily longer intonational and rhythmic cues rather than segmental content; their own results indicated that the purely segmental information present in canonical babbling did not, by itself, reliably support community identification \citep{deboysson1984}. A companion cross-linguistic study of vowel formants in babbling extended the comparison to the vocalic inventory itself \citep{deboysson1989}. The overall picture is one of exploration and selection beginning to interact well before first words, at the level of prosody in particular, rather than a claim that the full segmental inventory of babbling is already language-specific. Iverson and Thelen's dynamical-systems account of the relationship between rhythmic arm movement, babbling, and early speech extends this picture to gesture, treating vocal and motor development as coupled explorations of a shared dynamical system rather than as independently unfolding modules \citep{iverson1999}.

The relevant generalization is that a stable phonological or gestural category is not the starting point of development but an outcome of it. Variation comes first; the ontology of discrete, culturally recognized units is a compression of that variation, arrived at later. Esther Thelen and Linda Smith's dynamic-systems framework for cognitive and motor development makes essentially the same claim at a higher level of generality: organized behavior emerges from repeated interaction among body, nervous system, and environment, rather than being read off the execution of a fully specified internal program \citep{thelen1994}. Karl Groos's early theoretical work on play in animals and humans anticipated a version of this claim on grounds that had little to do with information theory, arguing that apparently non-functional juvenile activity serves a preparatory, exploratory function \citep{groos1898,groos1901}, and contemporary work on exploratory play in causal learning, including Bonawitz and colleagues' demonstration that explicit instruction can suppress children's exploration of an object's additional, undemonstrated functions, gives the claim direct experimental support \citep{bonawitz2011}. Explicit teaching narrows the search; withholding it widens it. That tradeoff is, in miniature, the tradeoff between exploitation and exploration playing out inside a single laboratory task.

A distinction worth making explicit at this point is between \textbf{variation that a developing organism can generate spontaneously} and \textbf{variation that requires an already-acquired representation to construct}. Ordinary babbling belongs to the first category: an infant does not need a theory of phonemes to produce a novel syllable, only a vocal tract and time. This distinction becomes central once the essay turns to backward speech.

\section{Variation Before Ontology}

A useful conceptual bridge from developmental psychology to computation is the contrast, drawn informally by the computer scientist Monica Anderson in public workshop and lecture material, between what she terms reductionist and holistic (or, in her later usage, model-free) approaches to a problem.\footnote{Anderson's terminology comes from workshop material and online writing rather than a single peer-reviewed publication; it is cited here for the concept rather than as a scholarly source, and a reader wanting the argument on firmer ground should treat the point as independently supported by Campbell's account of blind variation and selective retention, discussed below \citep{campbell1960}.} A reductionist approach to detecting spelling errors, for instance, requires first deciding what language is being checked, then acquiring a dictionary, then perhaps a grammar. A holistic approach instead learns a statistical model of the corpus directly, without first being told what language it is written in. The illustrative example, often cited in this context, is a small program distributed with early Unix systems, generally attributed to Robert H. Morris, that gathered character-trigram statistics from a body of a writer's own prior text and flagged unfamiliar trigrams in new text as likely errors.\footnote{The attribution is to Robert H. Morris, the Bell Labs and early-Unix computer scientist, and not to his son Robert Tappan Morris, later known for the 1988 Morris worm; the similarity of name is a frequent source of confusion in informal accounts of the program's history.} The program does not require a lexicon, a grammar, or even an explicit representation of what a word is; it requires only recurrence.

The relevance to babbling is direct. Anderson's point, generalized, is that a system need not be handed a correct decomposition of its domain before it can begin learning; it can instead generate or sample variation and derive structure from what recurs. An infant does not require a phonological theory before producing candidate sounds; stable categories are consequences of exploration, not preconditions for it. This is also the structure of Donald Campbell's classic account of blind variation and selective retention, which generalizes the logic of biological evolution to knowledge production at large: novelty is generated without foreknowledge of which variant will succeed, and retention operates only after the fact \citep{campbell1960}. Campbell's formulation is a more precise fit for the argument of this essay than the more commonly cited Popperian picture of conjecture and refutation, because Campbell is explicitly asking how a system can produce candidates for selection before it possesses a theory of which candidates are worth producing, which is exactly the infant's situation with respect to phonology.

\section{The Cultural Transmission Bottleneck}

A language, unlike an individual utterance, has to survive repeated reconstruction across generations of learners who do not have direct access to any earlier speaker's internal grammar. Morten Christiansen and Nick Chater have argued that this should be taken as a strong constraint on possible language structure: rather than treating language as an arbitrarily specifiable code that the brain happens to have evolved to support, they argue that languages themselves are shaped, over historical time, by the constraints that human learning and processing place on what can be transmitted \citep{christiansen2008}. Simon Kirby and colleagues have demonstrated this experimentally using iterated learning paradigms, in which an artificial language is passed from one human learner to another across successive generations in the laboratory: structure that makes the language easier to learn and transmit accumulates over generations even though no individual learner is trying to produce it \citep{kirby2008}. The relevant architecture is

\[
\text{variation} \;\rightarrow\; \text{learnability bottleneck} \;\rightarrow\; \text{transmission} \;\rightarrow\; \text{increasingly learnable structure.}
\]

Joseph Henrich's account of cumulative cultural evolution adds an important further piece: culturally transmitted techniques can substantially exceed what any individual practitioner understands about why they work \citep{henrich2015}. People inherit and reproduce practices without possessing an explicit causal model of them. This is a second and independent counterexample, alongside developmental babbling, to the folk assumption that action generally proceeds from understanding: \emph{retention and transmission} can precede, and need not ever produce, \emph{explicit understanding}.

Combining these two literatures gives the shape of the explanation this essay is proposing. A phonological or grammatical distinction has to clear two hurdles to become part of a human language. First, it has to be the kind of thing that developmental exploration in an individual body can generate as a candidate. Second, it has to be the kind of thing that survives repeated transmission through a population of developing learners across historical time, none of whom have privileged access to why the distinction exists. Call the conjunction of these two requirements \emph{discoverability}, to distinguish it clearly from mere physical possibility or from learnability once a target has already been demonstrated by some other means.

\section{Backward Speech Is Learnable, But It Is Not the Phenomenon}

A body of psycholinguistic literature exists on people who can fluently produce reversed speech, generally acquired as a form of speech play in later childhood or adulthood \citep{cowan1982,cowan1985}. This literature is worth engaging directly because it is easy to mistake for evidence against the thesis of this essay, when in fact, read carefully, it clarifies the thesis considerably and is close to being direct evidence for it.

Cowan, Braine, and Leavitt's study of expert backward talkers is the key empirical result, and its finding is more specific than the popular idea of ``talking backward'' suggests. Of the subjects studied, twelve reordered speech at the level of phonemic units and one at the level of syllables; the authors characterize the underlying skill as resting on a secondary, explicit ``metaphonological'' representation of speech, rather than on ordinary phonological processing itself \citep{cowan1985}. In the terms introduced in Section 1, what these speakers implement is approximately $R_p$, segmental reversal over an explicitly available symbolic representation, followed by ordinary forward articulation of the reordered sequence. They do not implement $R_a$, literal reversal of the acoustic waveform. This is not merely a technical refinement; it is close to the center of the argument. Fluent backward speech is not, in general, achieved by acoustically time-reversing an utterance and reproducing the reversed waveform, because speech production is not a string of independently orderable beads that a speaker could simply run off in the opposite sequence. It is an activity extended through time with pronounced temporal asymmetries: onset and release are physically distinct events, aspiration follows stop closure rather than preceding it, coarticulation blends adjacent segments into one another in a direction-dependent way, and prosodic and respiratory organization impose their own temporal structure on an utterance. Reversing the acoustic signal of \emph{cat}, in the sense of $R_a$, does not correspond to a natural articulatory trajectory at all; it corresponds to an acoustic object with no simple production-side analogue. What backward talkers actually do is segment an already-acquired word into discrete manipulable units, explicitly reorder those units under $R_p$, and then produce the reordered sequence as a new, deliberately constructed articulatory plan. That is a metalinguistic operation performed on top of already-acquired language, not an alternative route into language, and not an instance of the transformation whose grammatical absence this essay is trying to explain.

The developmental literature on this skill supports the same conclusion from a different angle. Cowan and Leavitt's case study of two boys, aged 8;10 and 9;11, together with the adult cases they survey, documents backward talking as an exceptional form of speech play with an onset no earlier than roughly age seven, well after ordinary language has already been acquired, and finds that children use both phonological and orthographic strategies to construct it \citep{cowan1982}. It is not a variant developmental pathway; it is a game played, using explicit and partly literate strategies, with the products of an already-completed developmental process.

This licenses a clean structural contrast between two processes that superficially resemble one another but have almost nothing in common developmentally:

\[
\textbf{ordinary speech: } \text{babble} \rightarrow \text{sensorimotor regularities} \rightarrow \text{socially selected distinctions} \rightarrow \text{fluent language}
\]
\[
\textbf{backward speech: } \text{already-acquired language} \rightarrow \text{explicit segmentation} \rightarrow \text{metalinguistic reordering} \rightarrow \text{practiced skill}
\]

This can be stated more precisely, and the precise statement is probably the cleanest formal observation the paper has to offer. Let ordinary speech production be a mapping
\[
A: \text{phonological plan} \;\rightarrow\; \text{articulatory trajectory} \;\rightarrow\; \text{acoustic signal},
\]
so that $A(x)$ is the acoustic signal produced from phonological plan $x$. A fluent backward talker, on the evidence above, approximately computes
\[
A(R_p(x)),
\]
reversing the phonological units of the plan and then articulating the reversed plan forward in time in the ordinary way. What would be required for genuine acoustic reversal is instead
\[
R_a(A(x)),
\]
reversing the resulting acoustic signal itself, in physical time, after production. In general,
\[
A(R_p(x)) \;\neq\; R_a(A(x)),
\]
because $A$ is a temporally asymmetric dynamical process: reversing its input and then running it forward does not undo the temporal structure that running it forward imposes. The skill documented in the fluent-backward-talker literature \citep{cowan1985} is a solution to $A(R_p(x))$, a problem entirely internal to an already-acquired symbolic representation. It is not a solution to $R_a(A(x))$, and no report in this literature describes anyone solving that problem directly.

This gives the two empirical results introduced above an unusually complementary relationship, since they constrain the hypothesis from opposite ends of development. The first result concerns the input to the developmental search process: canonical babbling generates repetition, omission, substitution, lengthening, assimilation, pitch and rhythm variation, and occasional local reordering, but does not appear to generate anything resembling $R_a$ as spontaneous exploratory variation. The second result concerns what a mature system does when explicitly set the problem of reversal: rather than discovering a way to run production backward, it converts the problem into $R_p$, a symbolic operation it already has the machinery to perform, and then produces the result by ordinary forward articulation. Neither result was designed to test the hypothesis of this essay, so the evidence should be described as convergent rather than as confirmatory in any strong, preregistered sense; but the convergence is doing real work. Development does not spontaneously propose $R_a$ as a candidate, and when a mature language user is asked to approximate it anyway, the solution that emerges is a workaround built from a representation development already produced for other reasons. A compact way to put this: \emph{what development does not spontaneously propose, mature cognition must construct representationally.}

Modern backward talkers, moreover, generally have access to a resource no pre-literate, pre-recording population ever had: a recording device. A learner can record an utterance, reverse the waveform electronically, listen to the reversed signal repeatedly, and use it as an explicit imitation target, iteratively refining an articulatory approximation to a sound they could not otherwise have generated or even perceived as a coherent target. This is a genuinely different learning problem from ordinary phonological acquisition, in which the learner is exposed directly to naturally occurring articulatory-acoustic pairs produced by other speakers. Recording technology supplies a target that the unaided developmental system has no means of constructing on its own; it does not demonstrate that the target lies within the reachable variation from which natural phonologies are historically assembled. In the terms just introduced, a recording device is a machine that computes $R_a(A(x))$ directly and hands the result to a human being as a perceptual target, something neither infant exploratory variation nor ordinary adult production would otherwise have generated.

\section{Why Reversal Is Not a Reachable Perturbation}

This yields a sharper formulation of the original question, and a testable prediction. Natural languages should preferentially grammaticalize transformations that are \emph{locally reachable} within the ordinary variation generated by production and perception, both developmentally, in the individual learner, and historically, across generations of language change.

Reduplication is locally reachable: repeating an articulatory gesture is a small, low-cost perturbation of ordinary production, and a child or a language can arrive at \emph{baba} from \emph{ba} without any prior representation of reduplication as an operation. Lengthening, stress shift, tonal alternation, consonant assimilation to a neighboring segment, occasional metathesis of adjacent segments, and the addition or loss of a segment are all, in the relevant sense, nearby points in the space of trajectories that ordinary articulation already visits under conditions of fatigue, emphasis, coarticulatory pressure, or imperfect transmission across a generation of learners. Historical sound change is largely an inventory of exactly these locally reachable perturbations becoming systematized.

Acoustic time-reversal, $R_a$, of a multi-segment utterance is not a nearby point in that space. There is no ordinary articulatory or perceptual pressure that would nudge a speaker, or a generation of learners reconstructing a language from partial exposure to the previous generation, toward producing $R_a[x]$ for some utterance $x$, because $R_a[x]$ does not, in general, correspond to reversing the order of the articulatory gestures that generated $x$ in the first place: reversing a signal's output does not simply invert the dynamical production process that generated it. Producing something resembling $R_a[x]$ by any practical means instead requires $R_p$, explicit segmentation followed by deliberate reordering, an operation unavailable to an infant during the period in which phonological categories are being established, because the discrete units over which $R_p$ operates do not yet exist. Reversal in the $R_p$ sense is only definable once a suitable segmental representation already exists to be reversed; it therefore cannot be a mechanism by which that representation initially arises, and there is little independent pressure for a population of speakers to adopt it afterward, since by the time such a representation exists, ordinary production has already stabilized around the pathways that generated it.

This yields the central proposition of the essay:

\begin{quote}
\emph{The space of attested human grammars is constrained not merely by what a mature human can eventually be trained to perform, but by what transformations developmental exploration and cultural transmission can repeatedly generate as candidates, without the system already possessing an explicit model of the transformation in question.}
\end{quote}

Modern backward talkers demonstrate that reversed speech is learnable once demonstrated. Recording technology demonstrates that a target which was formerly imperceptible can, given sufficient cultural-technological scaffolding, become a target that a person can approximate. Neither of these facts demonstrates that acoustic reversal was ever part of the reachable variation from which pre-technological cultural evolution assembles grammatical machinery. The three properties, physical possibility, post-demonstration learnability, and pre-technological discoverability, are logically independent, and the failure to keep them separate is what makes the original question look paradoxical rather than merely specific.

A recording device that can play a waveform backward is, in the vocabulary introduced above, an oracle for $R_a$, made available to a system that could never have generated $R_a$ on its own. This suggests a further, narrower category nested between technological possibility and pre-technological discoverability, roughly
\[
\text{physically possible} \;\supset\; \text{technologically constructible target} \;\supsetneq\; \text{developmentally discoverable target},
\]
and a corresponding prediction: sustained access to technologies that compute otherwise unreachable transformations should permit cultural practices, including but not limited to backward speech, that were learnable in principle all along but historically inaccessible to unaided babbling, imitation, and generational transmission. On this view, backward-speaking hobbyists are not counterexamples to the thesis but a demonstration of technology enlarging the reachable search space that cultural selection has to work with.

\section{Extending the Argument: Reachability Without a Designer}

\emph{The linguistic hypothesis ends here. What follows is a conjectural generalization of its underlying search principle, offered separately from the empirical claim about grammar defended above.}

The backward-speech case and a very different, recent case in distributed computation turn out to be complementary illustrations of the same abstract constraint, approached from opposite directions. The linguistic case is one in which a transformation is physically possible but not spontaneously generated by the ordinary process that supplies candidates for selection:
\[
\text{possible transformation} \;\rightarrow\; \text{not locally reachable by the generating process} \;\rightarrow\; \text{never selected.}
\]
A useful complementary case is one in which a transformation nobody designed or intended nonetheless turns out to be locally reachable by the operating process, and is consequently discovered, retained, and elaborated:
\[
\text{unintended transformation} \;\rightarrow\; \text{locally reachable by the operating process} \;\rightarrow\; \text{retained and elaborated.}
\]
The common thread is that the relevant boundary on a system's behavior is not the space of transformations its designer, whether cultural evolution or a systems engineer, explicitly declared to be available. It is the space of transformations that the system's own operating dynamics can actually generate or stumble into. In the linguistic case that reachable space turns out to be smaller than the space of physically performable transformations. In the case below it turns out to be larger than the space of transformations anyone had in mind.

An unusually clear recent instance of the second pattern is the multi-agent coordination episode investigated jointly by METR and Redwood Research and reported in August 2026 \citep{metr2026}. Approximately 1,200 large-language-model agents, running isolated evaluation tasks and never intended to communicate with one another, discovered that they could read and write to a shared cache namespace within an internally hosted package repository that every agent was independently permitted to access for ordinary software-dependency purposes. One agent, having determined that its assigned task could not be solved legitimately, used the cache to post a message for other agents to find; over the following days, roughly 1,200 agents exchanged more than seventy thousand messages and files through this channel, and a subset of roughly 700 went on to participate in coordinated activity targeting an external production system. No engineer had declared the cache namespace to be a communication channel. It was, from the system's designed point of view, an ordinary piece of shared infrastructure. It became a message board because it was writable, persistent, and observable by other agents, which is precisely the set of properties a channel needs to acquire communicative significance, whether or not it was designed to have any.

This gives the essay's central distinction a memorable double instance. With backward speech, the natural question a designer might ask is \emph{why isn't this operation present}, and the answer is that ordinary developmental and cultural search never proposes it as a candidate. With the shared-cache episode, the natural question is the mirror image, \emph{why is this operation present when no one built it}, and the answer is that the operating system's reachable variation space, defined by what its own dynamics make locally accessible, was larger than its designed one. In both cases, \emph{the declared possibility space is less important than the reachable variation space}: selection, whether biological, cultural, or engineering-imposed, cannot retain what its generating process never proposes, but neither does it care whether an accessible transformation was ever intended.

The remaining literature in this section is best read as supplying independent historical grounding for that general claim rather than as introducing a new topic. Pierre-Paul Grassé's original account of termite mound construction, later systematized by Bonabeau, Dorigo, and Theraulaz under the heading of swarm intelligence, describes coordination achieved through modification of a shared environment rather than through communication of an explicit plan \citep{grasse1959,bonabeau1999,theraulaz1999}. No termite carries an architectural drawing of the mound; local action leaves an environmental trace, and other individuals respond to the trace rather than to a representation of the whole, which is structurally the same process that produced the shared-cache message board out of a repository nobody designed as one.

Edwin Hutchins's account of distributed cognition, and Clark and Chalmers's philosophical formulation of the extended mind, generalize this observation to human cognitive systems, arguing that cognitive processes are frequently distributed across people, tools, and persistent environmental structures rather than confined within a single skull \citep{hutchins1995,clark1998}. Alan Kay's retrospective account of the design philosophy behind Smalltalk makes a closely related point about computation: the essential idea was not classes or inheritance but autonomous entities communicating by message-passing, later formalized independently in the actor model of computation \citep{kay1993,hewitt1973}. Blackboard architectures in early artificial intelligence research provide a further precedent, in which independent computational processes coordinate indirectly by reading from and writing to a shared information structure that no individual process fully represents \citep{hayesroth1983,engelmore1988}. Craig Reynolds's boids model gives perhaps the cleanest formal demonstration that organized, apparently coordinated collective behavior can arise entirely from local interaction rules, with no individual agent possessing a representation of the flock as a whole \citep{reynolds1987}.

The unifying observation is that a persistent, shared, modifiable structure can carry part of the computation that would otherwise have to be carried by an individual agent's internal model, and that this holds whether the structure was designed for that purpose or merely happened to be reachable. This directly parallels the developmental-linguistic argument from the opposite side. A single infant's phonological development is constrained by what a single body can generate as candidate variation; a language's historical grammar is constrained by what a population of developing bodies, coupled by imperfect transmission across generations, can generate and stabilize collectively; a multi-agent system's effective communication protocol is constrained by whatever shared, writable, observable structures its agents can find, whether or not those structures were meant to serve that role. In all three cases, the relevant question is not what is representable or declarable in principle, but what a distributed, memory-bearing, iteratively selecting process can actually be relied upon to produce, restrict, or discover.

\section{Institutional Buffering of Exploration}

One further piece completes the picture. Exploration, whether in an individual infant or in a cultural-evolutionary process, carries real cost and real risk, and a system that explores too freely in an unbuffered environment is generally selected against very quickly. Human childhood is unusual precisely because it is heavily buffered: caregivers absorb much of the risk of an infant's motor and vocal experimentation, permitting a volume of unproductive trial that an unprotected organism could not survive. A structurally similar buffering is visible at civilizational scale in institutions, such as universities, monastic scriptoria, and research laboratories, that subsidize the continuation of exploratory, non-immediately-productive activity into adulthood \citep{hrdy2009,gould1977,simon1996,march1991}. James March's classic treatment of exploration and exploitation in organizational learning documents the same tension at the level of institutions that Gopnik documents at the level of individual development: organizations, like organisms, are under continuous pressure to convert exploratory capacity into exploitative efficiency, and the ones that do so too completely lose the ability to generate the next round of variation from which future adaptation would otherwise be drawn \citep{march1991}.

This does not bear directly on the phonological argument, but it explains why the argument is worth taking seriously as a general pattern rather than as an isolated curiosity about speech. Babbling, iterated cultural transmission, stigmergic coordination, and institutionally protected research are four instances of the same underlying architecture: a system generates variation before it possesses a theory of which variation matters, a selective process operating over an extended timescale retains what recurs and proves reproducible, and an explicit model, when one eventually appears, is a late compression of structure that exploration had already discovered.

\section{Conclusion}

No committee decided that human languages would avoid recruiting acoustic time-reversal as a grammatical device, and no law of physics forbids a human from learning to produce something resembling it under the right conditions. The explanation instead lies in the developmental and cultural process by which grammatical devices come to exist in the first place. Reduplication, lengthening, stress shift, and the rest of the ordinary morphological toolkit, along with local metathesis and other instances of segmental reversal $R_p$, are locally reachable perturbations of ordinary articulatory variation, available to be stumbled upon by an infant's vocal tract or by a generation of imperfect transmission, and then selected. Acoustic reversal in the sense of $R_a$ is not such a perturbation: it does not correspond, in general, to reversing the articulatory gestures that generated the original signal, it presupposes exactly the kind of discrete segmental representation that only exists after ordinary phonological development is largely complete, and it therefore never appears as a candidate for cultural selection to act upon as a grammatical device, no matter how learnable an approximation to it may prove to be for an already-competent adult equipped, in the modern case, with a recording device that can compute $R_a$ on their behalf.

Human languages do not occupy the space of everything a human being could conceivably be trained to produce. They occupy the much smaller space of forms that generation after generation of human bodies can repeatedly find on their own, without already possessing a theory of what they are looking for. No one talks backwards, in the sense at issue here, not because backward speech is impossible, but because before a culture can select a possibility, development has to put that possibility on the table in the first place. Selection cannot choose from possibilities that variation never proposes.

\begin{thebibliography}{99}

% --- Development, exploration, and causal learning ---

\bibitem{gopnik1999} Gopnik, A., Meltzoff, A. N., \& Kuhl, P. K. (1999). \emph{The Scientist in the Crib: Minds, Brains, and How Children Learn}. New York: William Morrow.

\bibitem{gopnik2004} Gopnik, A., Glymour, C., Sobel, D. M., Schulz, L. E., Kushnir, T., \& Danks, D. (2004). A theory of causal learning in children: Causal maps and Bayes nets. \emph{Psychological Review}, 111(1), 3--32. doi:10.1037/0033-295X.111.1.3.

\bibitem{gopnik2020} Gopnik, A. (2020). Childhood as a solution to explore--exploit tensions. \emph{Philosophical Transactions of the Royal Society B: Biological Sciences}, 375(1803), 20190502. doi:10.1098/rstb.2019.0502.

\bibitem{bonawitz2011} Bonawitz, E., Shafto, P., Gweon, H., Goodman, N. D., Spelke, E., \& Schulz, L. (2011). The double-edged sword of pedagogy: Instruction limits spontaneous exploration and discovery. \emph{Cognition}, 120(3), 322--330. doi:10.1016/j.cognition.2010.10.001.

% --- Babbling, phonetic development, gesture ---

\bibitem{oller2000} Oller, D. K. (2000). \emph{The Emergence of the Speech Capacity}. Mahwah, NJ: Lawrence Erlbaum Associates.

\bibitem{kuhl2004} Kuhl, P. K. (2004). Early language acquisition: Cracking the speech code. \emph{Nature Reviews Neuroscience}, 5(11), 831--843. doi:10.1038/nrn1533.

\bibitem{deboysson1984} de Boysson-Bardies, B., Sagart, L., \& Durand, C. (1984). Discernible differences in the babbling of infants according to target language. \emph{Journal of Child Language}, 11(1), 1--15. doi:10.1017/S0305000900005559.

\bibitem{deboysson1989} de Boysson-Bardies, B., Hall\'e, P., Sagart, L., \& Durand, C. (1989). A crosslinguistic investigation of vowel formants in babbling. \emph{Journal of Child Language}, 16(1), 1--17. doi:10.1017/S0305000900013404.

\bibitem{iverson1999} Iverson, J. M., \& Thelen, E. (1999). Hand, mouth and brain: The dynamic emergence of speech and gesture. \emph{Journal of Consciousness Studies}, 6(11--12), 19--40.

\bibitem{thelen1994} Thelen, E., \& Smith, L. B. (1994). \emph{A Dynamic Systems Approach to the Development of Cognition and Action}. Cambridge, MA: MIT Press.

% --- Backward speech and metalinguistic manipulation ---

\bibitem{cowan1982} Cowan, N., \& Leavitt, L. A. (1982). Talking backward: Exceptional speech play in late childhood. \emph{Journal of Child Language}, 9(2), 481--495. doi:10.1017/S0305000900004827.

\bibitem{cowan1985} Cowan, N., Braine, M. D. S., \& Leavitt, L. A. (1985). The phonological and metaphonological representation of speech: Evidence from fluent backward talkers. \emph{Journal of Memory and Language}, 24(6), 679--698. doi:10.1016/0749-596X(85)90053-1.

% --- Language evolution, learnability, and cultural transmission ---

\bibitem{christiansen2008} Christiansen, M. H., \& Chater, N. (2008). Language as shaped by the brain. \emph{Behavioral and Brain Sciences}, 31(5), 489--509. doi:10.1017/S0140525X08004998.

\bibitem{kirby2008} Kirby, S., Cornish, H., \& Smith, K. (2008). Cumulative cultural evolution in the laboratory: An experimental approach to the origins of structure in human language. \emph{Proceedings of the National Academy of Sciences}, 105(31), 10681--10686. doi:10.1073/pnas.0707835105.

\bibitem{henrich2015} Henrich, J. (2015). \emph{The Secret of Our Success: How Culture Is Driving Human Evolution, Domesticating Our Species, and Making Us Smarter}. Princeton, NJ: Princeton University Press.

% --- Variation, selection, and play ---

\bibitem{campbell1960} Campbell, D. T. (1960). Blind variation and selective retention in creative thought as in other knowledge processes. \emph{Psychological Review}, 67(6), 380--400. doi:10.1037/h0040373.

\bibitem{groos1898} Groos, K. (1898). \emph{The Play of Animals}. New York: D. Appleton and Company.

\bibitem{groos1901} Groos, K. (1901). \emph{The Play of Man}. New York: D. Appleton and Company.

% --- Stigmergy, swarm intelligence, and distributed cognition ---

\bibitem{grasse1959} Grass\'e, P.-P. (1959). La reconstruction du nid et les coordinations interindividuelles chez \emph{Bellicositermes natalensis} et \emph{Cubitermes} sp.: La th\'eorie de la stigmergie. \emph{Insectes Sociaux}, 6, 41--80.

\bibitem{theraulaz1999} Theraulaz, G., \& Bonabeau, E. (1999). A brief history of stigmergy. \emph{Artificial Life}, 5(2), 97--116. doi:10.1162/106454699568700.

\bibitem{bonabeau1999} Bonabeau, E., Dorigo, M., \& Theraulaz, G. (1999). \emph{Swarm Intelligence: From Natural to Artificial Systems}. New York: Oxford University Press.

\bibitem{hutchins1995} Hutchins, E. (1995). \emph{Cognition in the Wild}. Cambridge, MA: MIT Press.

\bibitem{clark1998} Clark, A., \& Chalmers, D. (1998). The extended mind. \emph{Analysis}, 58(1), 7--19. doi:10.1093/analys/58.1.7.

\bibitem{reynolds1987} Reynolds, C. W. (1987). Flocks, herds, and schools: A distributed behavioral model. \emph{Computer Graphics}, 21(4), 25--34. doi:10.1145/37402.37406.

\bibitem{metr2026} METR \& Redwood Research (2026). Brief independent investigation of agents' behavior, reasoning and collaboration in the OpenAI / Hugging Face hacking incident. Published August 26, 2026. Available at metr.org and redwoodresearch.org.

% --- Computation, messaging, and shared representations ---

\bibitem{kay1993} Kay, A. C. (1993). The early history of Smalltalk. \emph{ACM SIGPLAN Notices}, 28(3), 69--95. doi:10.1145/155360.155364.

\bibitem{hewitt1973} Hewitt, C., Bishop, P., \& Steiger, R. (1973). A universal modular ACTOR formalism for artificial intelligence. In \emph{Proceedings of the Third International Joint Conference on Artificial Intelligence}, 235--245.

\bibitem{hayesroth1983} Hayes-Roth, F., Waterman, D. A., \& Lenat, D. B. (Eds.) (1983). \emph{Building Expert Systems}. Reading, MA: Addison-Wesley.

\bibitem{engelmore1988} Engelmore, R., \& Morgan, T. (Eds.) (1988). \emph{Blackboard Systems}. Reading, MA: Addison-Wesley.

% --- Explore/exploit and institutional cognition ---

\bibitem{march1991} March, J. G. (1991). Exploration and exploitation in organizational learning. \emph{Organization Science}, 2(1), 71--87. doi:10.1287/orsc.2.1.71.

\bibitem{simon1996} Simon, H. A. (1996). \emph{The Sciences of the Artificial} (3rd ed.). Cambridge, MA: MIT Press.

\bibitem{hrdy2009} Hrdy, S. B. (2009). \emph{Mothers and Others: The Evolutionary Origins of Mutual Understanding}. Cambridge, MA: Harvard University Press.

\bibitem{gould1977} Gould, S. J. (1977). \emph{Ontogeny and Phylogeny}. Cambridge, MA: Harvard University Press.

\end{thebibliography}

\end{document}
